\begin{lemma}
\label{editorial-source-lemma-make-geometrically-irreducible}
Let $K/k$ be a field extension. Consider the subextension $K/k'/k$ consisting
of elements separably algebraic over $k$. Then $K$ is geometrically irreducible
over $k'$. If $K/k$ is a finitely generated field extension, then
$[k' : k] < \infty$.
\end{lemma}

\begin{proof}
The original set $k'$ is a subfield
by Fields, Lemma \ref{fields-lemma-separable-elements}.
If $a\in K$ is separably algebraic
over $k'$, transitivity of separable
algebraic extensions (Fields, Lemma
\ref{fields-lemma-separable-permanence})
makes $a$ separably algebraic over
$k$. Hence $a\in k'$ by its
definition. Lemma
\ref{editorial-source-lemma-geometrically-irreducible-separable-elements}
now proves geometric irreducibility
of the original $K/k'$.

For completeness, retain the full
degree argument of Fields, Lemma
\ref{fields-lemma-algebraic-closure-in-finitely-generated}.
Choose a finite transcendence basis
$x_1,\ldots,x_r$ of the finitely
generated extension $K/k$ and put
$$
F=k(x_1,\ldots,x_r),\qquad n=[K:F]<\infty.
$$
Here $r=0$ is allowed, with $F=k$.
The remaining members of a finite
generating set for $K/k$ are
algebraic over $F$, so adjoining
them successively proves the
displayed finiteness. Let $E/k$
be any original finite algebraic
subextension of $K/k$, and choose
a basis $e_1=1,e_2,\ldots,e_d$ of
$E/k$, where $d=[E:k]$.

The original $x_i$ remain
algebraically independent over $E$.
Here is a direct proof retaining
every coefficient. For nonzero
$P\in E[T_1,\ldots,T_r]$, let $M_P$
be the matrix of multiplication
by $P$ on the free $k[T_1,\ldots,T_r]$-module
$E[T_1,\ldots,T_r]$ in the basis
$e_1,\ldots,e_d$. This multiplication
is injective because the latter
polynomial ring is a domain.
After localization to $k(T_1,\ldots,T_r)$,
the matrix is an injective square
linear map on a finite-dimensional
vector space, hence
$N_P=\det(M_P)\neq0$ in $k[T_1,\ldots,T_r]$.
Write $H_P\in E[T_1,\ldots,T_r]$
for the element whose coordinate
column is $\operatorname{adj}(M_P)(1,0,\ldots,0)^t$.
The full adjugate identity gives
$$
P H_P=N_P.
$$
If $P(x_1,\ldots,x_r)=0$, this
identity would imply
$N_P(x_1,\ldots,x_r)=0$, contradicting
the original independence over $k$.
This proves the asserted independence
over $E$, also for $r=0$.

Consequently the original evaluation
map induces inverse isomorphisms
$$
E\otimes_k F\longrightarrow E(x_1,\ldots,x_r),
\qquad a\otimes h(x)\longmapsto a h(x).
$$
To specify the inverse, write a
numerator in the basis $e_i$ and
retain each original denominator
$Q(x)$ with $Q\in E[T_1,\ldots,T_r]\setminus\{0\}$.
Use exactly the identity
$Q^{-1}=H_Q/N_Q$ just proved,
then expand the numerator $P H_Q$
in the same basis. This gives
its unique expression
$\sum_i e_i f_i(x)$ with $f_i\in F$,
and the inverse sends it to
$\sum_i e_i\otimes f_i(x)$.
Uniqueness follows by multiplying
all denominators in $k[T_1,\ldots,T_r]$
and comparing polynomial coefficients
in the basis $e_i$; no denominator
vanishes at the original $x$.
Thus these are well-defined inverse
maps and
$$
[E:k]=[E(x_1,\ldots,x_r):F]\leq[K:F]=n.
$$

In fact this bounds the full
algebraic constant field $k_a$
inside $K$, not only $k'$.
The degrees of its finite
subextensions lie among the
integers $1,\ldots,n$, so choose
one, $E_0$, of maximal degree.
For any $a\in k_a$, the original
$E_0(a)/k$ is finite and has
degree at least $[E_0:k]$.
Maximality and degree multiplicativity
give $[E_0(a):E_0]=1$. Hence every
$a$ lies in $E_0$, proving
$k_a=E_0$ and the precise bound
$$
[k':k]\leq[k_a:k]\leq[K:k(x_1,\ldots,x_r)].
$$

There is also an exact description
of the algebraic intermediate
fields $k\subseteq L\subseteq K$
over which $K$ is geometrically
irreducible: they are exactly
those containing $k'$.
For necessity, every $a\in k'$
is separably algebraic over $L$,
because its minimal polynomial
over $L$ divides its original
separable polynomial over $k$.
Geometric irreducibility of $K/L$
and the preceding criterion force
$a\in L$. For sufficiency, suppose
$k'\subseteq L$ and $L/k$ algebraic.
The full algebraic constant field
$k_a/k'$ is purely inseparable
by Fields, Lemma
\ref{fields-lemma-separable-first}.
If $a\in K$ is separably algebraic
over $L$, it is algebraic over
$k$, hence belongs to $k_a$.
Thus it is also purely inseparable
over $L$: some original $p$-power
lies in $k'\subseteq L$ in positive
characteristic, and in characteristic
zero $k_a=k'$. An element both
separable and purely inseparable
has degree one, since its minimal
polynomial is separable and divides
a polynomial with only one distinct
root. Therefore $a\in L$, and the
criterion proves sufficiency.
The restriction $L/k$ algebraic
is part of this characterization;
it is not removed.
\end{proof}